% Lecture 05: Классы и инкапсуляция

\section{Лекция 5. Классы и инкапсуляция}
\label{sec:lecture-05}

В блоке~4 большой \texttt{main} утилиты инвентаризации был разделён на
функции. Но строка отчёта всё ещё представлена \texttt{BlockInfo} с открытыми
полями. Имя, пути и результаты проверок постоянно путешествуют вместе. Caller
может записать \texttt{lecture\_exists=true}, оставить нулевое число строк или
связать имя одного блока с путём другого. Типы полей корректны, состояние в
целом --- нет.

\begin{importantbox}[Главный вопрос]
Как создать тип предметной области, который хранит связанные данные рядом,
предлагает только осмысленные операции и сокращает число мест, где можно
нарушить согласованность?
\end{importantbox}

\subsection{Зачем программе собственный тип}

В языке уже есть \texttt{int}, \texttt{bool} и pointer. Библиотека даёт
строки, контейнеры и пути: \texttt{std::string}, \texttt{std::vector},
\texttt{std::filesystem::path}. Но задача говорит не о случайном числе и
строке, а о \emph{блоке курса}. Набор переменных плохо выражает сущность:

\begin{cppcode}
int number{5};
std::string slug{"classes_and_encapsulation"};
bool lecture_exists{true};
bool examples_exist{true};
\end{cppcode}

Если параметры почти всегда передаются вместе, это сигнал проверить, не
являются ли они одним понятием. Не всякая длинная parameter list требует
класса, но повторяющаяся группа данных и операций часто обнаруживает
потерянный тип. Класс позволяет создать новый пользовательский тип, соединяя
состояние с операциями над ним.

\begin{definition}
\textbf{Class} задаёт пользовательский тип. \textbf{Object}, или
\textbf{instance}, --- конкретный объект этого типа со своим состоянием и
lifetime. Class --- описание; \texttt{first} и \texttt{second} ниже --- два
разных объекта.
\end{definition}

\begin{cppcode}
struct BlockLabel {
    int number;
    std::string slug;
};

BlockLabel first{4, "functions_references_and_const"};
BlockLabel second{5, "classes_and_encapsulation"};
\end{cppcode}

\texttt{number} и \texttt{slug} --- data members, то есть объекты, являющиеся
частями каждого \texttt{BlockLabel}. У двух instances свои members. Точка
выбирает member конкретного объекта. Поэтому \texttt{first.number} не
обращается к состоянию \texttt{second}.

\subsection{Struct и class: языковое различие}

Обе формы могут содержать данные и функции:

\begin{cppcode}
struct PublicLabel {
    int number; // public по умолчанию
    std::string slug;
    bool valid() const { return number >= 1 && number <= 15; }
};

class HiddenLabel {
    int number_; // private по умолчанию
    std::string slug_;
};
\end{cppcode}

Ключевое различие здесь --- default access: members \texttt{struct} по
умолчанию \texttt{public}, members \texttt{class} --- \texttt{private}.
Обе формы способны иметь data members, member functions и access sections.
Миф «struct только для данных, class только для поведения» не правило C++.

Практическая эвристика: \texttt{struct} удобен для простой агрегатоподобной
записи с намеренно открытым представлением; \texttt{class} --- когда нужно
скрыть representation или защищать invariant. Это style guideline, не
ограничение языка.

\subsection{Member function и текущий объект}

Функция внутри определения класса называется member function (часто method):

\begin{cppcode}
class Progress {
public:
    bool is_complete() const {
        return lecture_exists_ && examples_exist_;
    }

private:
    bool lecture_exists_{};
    bool examples_exist_{};
};
\end{cppcode}

Call \texttt{first.is\_complete()} выбирает method через \texttt{.}. Он
работает с состоянием именно \texttt{first}. Интуитивная модель: у non-static
member function есть неявный object parameter. Внутри неё \texttt{this} ---
pointer на текущий объект; поэтому \texttt{lecture\_exists\_} означает member
текущего объекта и при необходимости записывается
\texttt{this->lecture\_exists\_}. Обычно явный \texttt{this->} здесь только
зашумляет код.

Это не означает копирование объекта при каждом call. Способ передачи скрытого
адреса зависит от реализации и ABI; языковая семантика задаёт, к какому
объекту относится вызов. Method принимает обычные parameters, возвращает
values и может менять текущий non-const object.

Compiler проверяет типы arguments, доступность method и допустимость операции
для выбранного объекта. Runtime выполняется обычное тело функции и меняется
member конкретного instance.

\subsection{Definition вне class и operator ::}

Короткие methods можно определить внутри class definition. Длинное тело часто
выносят после объявления:

\begin{cppcode}
class Counter {
public:
    void add(int amount);
    int value() const;

private:
    int value_{};
};

void Counter::add(int amount) {
    value_ += amount;
}

int Counter::value() const {
    return value_;
}
\end{cppcode}

\texttt{Counter::add} означает: имя \texttt{add} находится в scope класса
\texttt{Counter}. Это тот же scope resolution operator \texttt{::}, что в
\texttt{std::cout}, но слева class, а не namespace. Declaration и definition
должны согласоваться, включая parameters, return type и завершающий
\texttt{const}.

Definition внутри класса удобно для маленькой операции и по правилам языка
--- inline definition; это не обещание машинного inlining. Вынос тела
помогает отделить interface от implementation.

\subsection{Public, private и граница доступа}

\texttt{public} members образуют интерфейс, доступный caller. \texttt{private}
members доступны самому классу (и friends, которые сейчас не нужны).

\begin{cppcode}
class CourseProgress {
public:
    bool is_complete() const { return lecture_ && examples_; }
    void mark_lecture_found() { lecture_ = true; }

private:
    bool lecture_{};
    bool examples_{};
};

CourseProgress progress;
progress.mark_lecture_found(); // допустимо
// progress.lecture_ = false;  // compile-time error: private
\end{cppcode}

Access control --- compile-time mechanism. Корректная программа не получает
runtime-проверку на каждое обращение к private field. Если caller называет
недоступный member, программа ill-formed, и compiler выдаёт diagnostic.

\begin{warningbox}[Private не защищает от атак]
\texttt{private} организует C++ interface, но не шифрует память и не создаёт
security boundary. UB, ошибочный raw pointer или debugger не останавливаются
словом \texttt{private}.
\end{warningbox}

Когда всё public, любое место программы может менять каждую деталь и зависит
от её формы. Смена представления требует искать callers, а корректность
поддерживается распределённой дисциплиной.

\subsection{Encapsulation как технический приём}

\begin{definition}
\textbf{Encapsulation} здесь означает совместное размещение связанного
состояния и операций, скрытие деталей representation и ограничение набора
допустимых операций публичным interface.
\end{definition}

Это не лозунг «положить поля под \texttt{private}». Полезный результат:
меньше раздельных параметров и комбинаций несогласованных данных; проверки и переходы состояния сосредоточены в небольшом числе methods; caller зависит от смысла операций, а не от способа хранения; representation можно менять без переписывания caller, пока сохраняется контракт interface.

Например, реализация может заменить два bool members на другое внутреннее
представление, оставив \texttt{is\_complete()} неизменным. Это возможность
design, а не автоматическое обещание binary compatibility: изменение layout
может быть существенно для отдельно собранных components.

\subsection{Invariant: что должно оставаться истинным}

\begin{definition}
\textbf{Invariant объекта} --- условие, которое должно оставаться истинным для
корректно используемого объекта на границах его публичных операций.
\end{definition}

Для \texttt{CourseBlock}: number находится в диапазоне 1--15, slug не пуст,
а lecture/examples paths согласованы с number и slug. Плохой public record
разрешает caller стереть slug, записать 99 и оставить старые paths. Хороший
design не предоставляет таких операций: paths вычисляются вместе с identity,
а public interface только читает измерения.

\texttt{private} сам по себе invariant не гарантирует. Method может записать
99 или ошибочно вычислить path. Гарантию создаёт совокупность interface и
implementation, подтверждённая тестами. Compiler знает типы и access rules,
но не выводит из предметной области диапазон 1--15.

В практике применяется factory-like operation \texttt{CourseBlock::inspect}.
Она строит согласованное состояние и возвращает его как value. Для создания
внутри используется минимальный default constructor; полная модель
constructors/destructors и более сильная инициализация --- тема блока~6.

\subsection{Const member functions}

\begin{cppcode}
class CourseBlock {
public:
    int number() const { return number_; }
    bool is_complete() const {
        return lecture_exists_ && examples_exist_;
    }

private:
    int number_{};
    bool lecture_exists_{};
    bool examples_exist_{};
};
\end{cppcode}

\texttt{const} после parameter list относится к object parameter, а не к
return type. Такой method обещает read-only access к обычному observable
member state: через него нельзя присвоить значение обычному non-mutable
member. Тему \texttt{mutable} сейчас не вводим.

\begin{cppcode}
void print_one(const CourseBlock& block) {
    std::cout << block.number() << ' ' << block.is_complete() << '\n';
}
\end{cppcode}

Связь с блоком~4 прямая: \texttt{const CourseBlock\&} даёт read-only alias без
копии, а вызываемые через него methods должны быть const. Попытка вызвать
modifying non-const method через const object --- compile-time error. Поэтому
const methods формируют читающую часть interface.

Const member function не обещает глобальную неизменность объекта: другой
non-const alias может его изменить. Она также не решает lifetime и ownership.

\subsection{Getter и setter: не цель инкапсуляции}

Getter бывает уместен: \texttt{number()} сообщает устойчивое свойство,
\texttt{lecture\_path()} предоставляет read-only view. Setter тоже нормален,
если независимое присваивание действительно допустимо
предметной области.

Но \texttt{getX()/setX()} для каждого private field часто лишь имитирует
public fields. Caller всё ещё создаёт бессмысленные переходы. Осмысленный
interface говорит \texttt{is\_complete()},
\texttt{mark\_lecture\_found()}, \texttt{missing\_parts()} или
\texttt{lecture\_path()}; названия выражают задачу и дают реализации место
поддержать invariant.

Нельзя оценивать encapsulation числом getters. Важнее, какие неправильные
действия interface сделал невозможными и какие изменения representation
перестали затрагивать callers.

\subsection{Responsibility, cohesion и free functions}

У маленького класса должна быть одна внятная responsibility.
\texttt{CourseBlock} представляет и инспектирует один блок. Он не разбирает
command line и не печатает общую таблицу.

\begin{definition}
На начальном уровне высокая \textbf{cohesion} означает, что members и methods
класса связаны одной ролью и его состоянием.
\end{definition}

Method естественен, когда операция требует состояния/invariant конкретного
объекта: \texttt{block.is\_complete()}. Операция, не использующая это
состояние, не обязана быть method. \texttt{print\_report(blocks)} и
\texttt{read\_arguments(...)} остаются free functions.

Заталкивание всех функций в один class создаёт God object: много причин для
изменения, сильную связанность и трудный interface. Плохой class design может
быть хуже ясных данных и нескольких функций.

\subsection{Composition: отношение has-a}

\begin{cppcode}
class CourseBlock {
    std::string slug_;
    std::filesystem::path lecture_path_;
    std::filesystem::path examples_path_;
};
\end{cppcode}

\texttt{CourseBlock} has-a string и два path objects. Их lifetimes являются
частью lifetime содержащего объекта. Уже знакомый
\texttt{std::vector<CourseBlock>} содержит элементы. Composition естественна,
когда модель действительно «имеет» часть; inheritance и virtual здесь не
нужны и появятся позже.

\subsection{Memory, lifetime и layout}

Объект class type имеет lifetime по тем же общим правилам. Локальный
\texttt{CourseBlock} живёт до выхода из scope; элемент vector --- пока
существует элемент и операция контейнера его не уничтожила. Pointer/reference
не продлевает lifetime.

Non-static data members являются частями complete object. Их создание и
уничтожение связано с содержащим объектом, но подробный порядок и управление
ресурсами относятся к блоку~6. \texttt{sizeof(CourseBlock)} можно измерить:
результат включает представления members и возможный padding. Нельзя по одному
измерению объявлять точный layout гарантией всех compiler/ABI.

Access specifiers не обязаны добавлять «флаг private» в object. Стандарт
задаёт layout ограниченно; alignment, padding и представления библиотечных
members зависят от реализации.

\subsection{Cost model без объектной магии}

Обычная non-virtual member function не требует отдельной объектной VM. Access
control проверяется при трансляции; \texttt{public}/\texttt{private} сами по
себе не runtime checks. Method может быть оптимизирован/inlined по тем же
общим причинам, что free function.

Нельзя утверждать, что method всегда быстрее или медленнее free function.
Значимая стоимость здесь приходит от filesystem calls, allocations string и
vector, копий и I/O. Encapsulation может не добавлять отдельной runtime платы,
но плохой API способен вызвать лишние копии или coupling. Zero-overhead intent
C++ --- возможность избежать обязательной доплаты, не обещание нулевой цены
любого design. Virtual dispatch и vtable относятся только к блоку~7.

\subsection{Эксперимент 1: реальный class-based refactoring}

Каталог
\texttt{examples/05\_classes\_and\_encapsulation/01\_course\_inventory\_classes}
развивает \texttt{course\_inventory\_refactored} блока~4. CLI, столбцы,
totals и классы exit codes сохранены.

В procedural версии \texttt{BlockInfo} открывал каждую деталь. Теперь private
members одного \texttt{CourseBlock} держат number, slug, согласованные paths
и измерения. Static public operation \texttt{inspect} создаёт их вместе.
Read-only methods \texttt{lecture\_lines()},
\texttt{example\_file\_count()} и \texttt{is\_complete()} являются const.
Setter для каждого field отсутствует: caller не должен переписывать измеренное
состояние.

\texttt{collect\_blocks}, \texttt{read\_arguments} и
\texttt{print\_report} остались free functions: они управляют коллекцией, CLI
и выводом, а не invariant одного блока. Поэтому \texttt{main} читает как
сценарий, а не как God object.

\begin{shellcode}
cd examples/05_classes_and_encapsulation/01_course_inventory_classes
cmake -S . -B build
cmake --build build
./build/course_inventory_classes --root ../../.. --block 5
ctest --test-dir build --output-on-failure
\end{shellcode}

Эквивалентность проверяется запуском старого и нового executable с одинаковыми
arguments, заменой только process-specific printed addresses и сравнением
stdout через \texttt{diff}. PASS: normalized output и exit codes совпали,
warnings отсутствуют, три CTest прошли.

\subsection{Эксперимент 2: encapsulation\_break\_lab}

В \texttt{02\_encapsulation\_break\_lab} public
\texttt{struct PublicBlock} принимает number 99, пустой slug и несогласованные
flags. \texttt{CourseProgress} скрывает state и предлагает
\texttt{for\_block}, \texttt{mark\_lecture\_found},
\texttt{mark\_examples\_found}, \texttt{is\_valid} и
\texttt{is\_complete}. Два instances изменяются независимо.

Free function \texttt{status\_line(const CourseProgress\&)} форматирует отчёт
через read-only interface: форматирование не обязано быть method. Const
reference вызывает const methods. Positive executable заканчивает
\texttt{ENCAPSULATION\_BREAK\_LAB=PASS}.

\begin{shellcode}
cd examples/05_classes_and_encapsulation/02_encapsulation_break_lab
cmake -S . -B build
cmake --build build
./build/encapsulation_break_lab
ctest --test-dir build --output-on-failure
\end{shellcode}

CTest вызывает настоящий compiler для двух ill-formed files.
\texttt{private\_access.cpp} пытается записать private \texttt{number\_};
\texttt{const\_mutation.cpp} вызывает non-const modifying method через const
object. PASS требует ненулевой exit status и непустую diagnostic в build.
Текст ошибки не подделывается и различается между GCC/Clang.

Sanitizers здесь N/A: проверяются compile-time запреты, а positive cases не
выполняют опасной memory operation. Намеренное UB не проверило бы access
control лучше.

\subsection{Сильные стороны и цена}

C++ позволяет создавать доменные типы, compiler проверяет access control и
const object contract, а ordinary encapsulation может не иметь отдельной
runtime платы. Programmer выбирает method для операции над состоянием и free
function для внешней операции.

Цена: бессмысленные getters/setters не создают хороший interface; compiler не
знает business invariant; слишком большой class усиливает coupling; lifetime
и ownership не исчезают. Constructors/destructors следующего блока нужны для
более сильных invariants и resource-owning types.

\subsection{Проверка понимания}

После блока нужно уметь:
1) различить class, object/instance, data member и member function; 2) назвать default-access различие \texttt{struct}/\texttt{class}; 3) объяснить operator \texttt{.}, текущий объект и базовую роль \texttt{this}; 4) описать public interface, private representation и границы защиты; 5) определить invariant и объяснить, почему private недостаточно; 6) прочитать definition method через \texttt{::}; 7) объяснить \texttt{const} после parameter list; 8) отличить domain operation от механического setter; 9) решить, должна ли операция быть method или free function; 10) описать composition, lifetime members и осторожный cost model; 11) показать, чем новый \texttt{CourseBlock} лучше procedural record; 12) интерпретировать обе реальные negative diagnostics.

\subsection{Источники для сверки}

Class definitions, members, access control, \texttt{this}, cv-qualification
member functions, member access и object lifetime сверены с C++17 working
draft N4659. Последовательность user-defined datatype, instance, fields,
methods, composition и \texttt{::} сопоставлена с MIT 6.096 Lecture~6; из
Lecture~7 использована только encapsulation, без inheritance/polymorphism.
Старые публичные учебные classes не приняты как современное правило design.
Style и target-based сборка сверены с локальными \textit{Modern C++ Tutorial}
и \textit{Modern CMake}.
